\documentclass{article}
\usepackage{amsmath, amssymb, amsthm}
\usepackage{fontspec}
\usepackage{etoolbox}
\AtBeginEnvironment{align}{\setcounter{equation}{0}}

\setmainfont{Kinnari}
[
	Extension = .ttf,
	UprightFont = *,
	BoldFont = *-Bold,
	ItalicFont = *-Italic,
	BoldItalicFont = *-BoldItalic
]

\title{Homework 20}
\author{68011008 Waroon Ragwongsiri}
\date{19/02/2026}

\begin{document}

\maketitle

% Q1
{\small
ณ ศาลแห่งหนึ่งมีการเบิกความดังนี้ \\
ก. ถ้าลูกความมีความผิด มีดต้องอยู่ในลิ้นชัก \\
ข. มีดไม่ได้อยู่ในลิ้นชัก หรือ ช๊อคซิปเห็นมีด \\
ค. หากไม่มีมีดอยู่ในวันปีใหม่ ช๊อคชิปก็จะไม่เห็นมีด \\
ง. ถ้ามีมีดอยู่ที่นั่นในวันปีใหม่ แสดงว่า มีดจะอยู่ในลิ้นชัก พร้อมๆ กับ ค้อนต้องอยู่ในโรงนาด้วย \\
จ. แต่เราทุกคนรู้ดีว่าค้อนไม่ได้อยู่ในโรงนา \\
จงหาข้อสรุปว่าลูกความมีความผิดหรือไม่ โดยใช้กระบวนการ direct deduction \\
let ลูกความมีความผิด = p \\
let มีดต้องอยู่ในลิ้นชัก = q \\
let ช๊อคซิปเห็นมีด = r \\
let มีมีดอยู่ในวันปีใหม่ = t \\
let ค้อนต้องอยู่ในโรงนา = s
}
\begin{align}
	&p \implies q \quad & Premise \\
	&\neg q \lor r \quad & Premise \\
	&\neg t \implies \neg q \quad & Premise \\
	&t \implies (q \land s) \quad & Premise \\
	&\neg s \quad & Premise \\
	&q \implies t \quad & 3 \ Contrapositive \\
	&p \implies t \quad & 1,6 \ Hypothetical Syllogism \\
	&p \implies (q \land s) \quad & 4,7 \ Hypothetical Syllogism \\
	&\neg p \lor (q \land s) \quad & 8 \ Material Implication \\
	&(\neg p \lor q) \land (\neg p \lor s) \quad & 9 \ Distribution \\
	&\neg p \lor s \quad & 10 \ Simplication \\
	&p \implies s \quad & 11 \ Material Implication \\
	&\neg p \quad & 5,12 \ Modus Tollens
\end{align}
\small
{
\therefore {ลูกความไม่มีความผิด}
}

\end{document}